% ============================================================
% 02-bezier-sympy-output-explained.tex
% Compile with:  pdflatex -shell-escape 02-bezier-sympy-output-explained.tex
% Run twice for table of contents.
% ============================================================

\documentclass[12pt, a4paper]{article}

% ---- Encoding & fonts ----------------------------------------
\usepackage[T1]{fontenc}
\usepackage[utf8]{inputenc}
\usepackage{lmodern}

% ---- Page geometry -------------------------------------------
\usepackage[
  top=2.5cm, bottom=2.5cm,
  left=2.8cm, right=2.8cm
]{geometry}

% ---- Colours -------------------------------------------------
\usepackage[dvipsnames]{xcolor}
\definecolor{codebg}{RGB}{248,248,248}
\definecolor{rulecolor}{RGB}{180,180,180}
\definecolor{examplebg}{RGB}{240,248,255}
\definecolor{examplerule}{RGB}{100,149,237}
\definecolor{notebg}{RGB}{255,250,230}
\definecolor{noterule}{RGB}{210,180,50}
\definecolor{diffbg}{RGB}{240,255,240}
\definecolor{diffrule}{RGB}{80,160,80}

% ---- Minted --------------------------------------------------
\usepackage{minted}
\setminted[text]{
  bgcolor    = codebg,
  frame      = leftline,
  framerule  = 1.5pt,
  rulecolor  = rulecolor,
  fontsize   = \small,
  breaklines = true,
}

% ---- Math ----------------------------------------------------
\usepackage{amsmath}
\usepackage{amssymb}

% ---- Boxes ---------------------------------------------------
\usepackage{mdframed}

\newmdenv[
  backgroundcolor  = examplebg,
  linecolor        = examplerule,
  linewidth        = 1.5pt,
  leftline=true, rightline=false, topline=false, bottomline=false,
  innerleftmargin=10pt, innerrightmargin=8pt,
  innertopmargin=6pt, innerbottommargin=6pt,
  skipabove=6pt, skipbelow=4pt,
]{examplebox}

\newmdenv[
  backgroundcolor  = notebg,
  linecolor        = noterule,
  linewidth        = 1.5pt,
  leftline=true, rightline=false, topline=false, bottomline=false,
  innerleftmargin=10pt, innerrightmargin=8pt,
  innertopmargin=6pt, innerbottommargin=6pt,
  skipabove=6pt, skipbelow=4pt,
]{notebox}

\newmdenv[
  backgroundcolor  = diffbg,
  linecolor        = diffrule,
  linewidth        = 1.5pt,
  leftline=true, rightline=false, topline=false, bottomline=false,
  innerleftmargin=10pt, innerrightmargin=8pt,
  innertopmargin=6pt, innerbottommargin=6pt,
  skipabove=6pt, skipbelow=4pt,
]{diffbox}

% ---- Hyperlinks ----------------------------------------------
\usepackage[colorlinks=true, linkcolor=black, urlcolor=MidnightBlue]{hyperref}

% ---- Misc ----------------------------------------------------
\usepackage{booktabs}
\usepackage{needspace}
\usepackage{parskip}
\usepackage{titlesec}
\usepackage{fancyhdr}
\usepackage{datetime2}

% ---- Section formatting --------------------------------------
\titleformat{\section}
  {\large\bfseries\color{MidnightBlue}}
  {\thesection.}{0.6em}{}[\vspace{-4pt}\rule{\linewidth}{0.4pt}]
\titleformat{\subsection}
  {\normalsize\bfseries\color{BrickRed}}
  {\thesubsection.}{0.5em}{}

% ---- Header / footer -----------------------------------------
\pagestyle{fancy}
\fancyhf{}
\rhead{\textcolor{gray}{\small 02-bezier-sympy-output-explained}}
\lhead{\textcolor{gray}{\small B\'ezier Symbolic --- Output Explanation}}
\cfoot{\thepage}
\renewcommand{\headrulewidth}{0.3pt}

% ---- Helpers -------------------------------------------------
\newcommand{\cd}[1]{\texttt{#1}}
\newcommand{\NEW}{\textcolor{diffrule}{\textbf{[New]}}\;}

% ---- Title ---------------------------------------------------
\title{%
  \vspace{1cm}
  {\LARGE\bfseries B\'ezier Curve --- Algebraic Expansion (Symbolic)}\\[0.4em]
  {\large Explanation of the Program Output}\\[0.4em]
  {\normalsize Case 2: Quartic 3-D, 5 Control Points}\\[0.4em]
  {\normalsize\texttt{02-bezier-sympy.jl}}
}
\author{E.R.\ Nurwijayadi}
\date{\DTMtoday}


% =============================================================
\begin{document}
% =============================================================

\maketitle
\thispagestyle{empty}

\begin{center}
\textit{Directed and curated by E.R.\ Nurwijayadi.
Code and explanations AI-generated.}
\end{center}

\begin{center}
\begin{minipage}{0.85\textwidth}
\small
This document explains the output of \cd{02-bezier-sympy.jl}
when run with \cd{CASE = 2}.
It is a companion to the output explanation for \cd{01-bezier-algebra.jl}.
The curve and control points are identical --- what changed is the
computation engine.
Script~01 used a hand-written \cd{Poly} type with floating-point arithmetic.
Script~02 uses \textbf{Symbolics.jl}, a computer algebra system,
keeping all values as exact fractions throughout.
\end{minipage}
\end{center}

\vspace{1em}\hrule\vspace{1.5em}
\tableofcontents
\newpage


% =============================================================
\needspace{8\baselineskip}
\section{What Changed from Script 01}
% =============================================================

\textit{Same curve. Same case. Same five steps.
Different engine under the hood ---
and the difference shows up in every line of the output.}

Script~01 (\cd{01-bezier-algebra.jl}) built its own \cd{Poly} type
and stored polynomial coefficients as \cd{Float64} numbers.
All arithmetic was floating-point.

Script~02 (\cd{02-bezier-sympy.jl}) replaces the hand-rolled
\cd{Poly} type entirely with \textbf{Symbolics.jl} --- Julia's
computer algebra system (CAS).
Instead of numbers, the program manipulates \emph{symbolic expressions}
in the variable \cd{u}.
Coefficients are stored as \cd{Rational\{Int\}} --- exact integer
fractions --- so there is no floating-point rounding anywhere
until the very last step when a decimal is explicitly requested.

\begin{diffbox}
\textbf{The three visible differences in the output:}

\medskip
\begin{tabular}{lp{4.5cm}p{4.5cm}}
\toprule
Location & Script 01 & Script 02 \\
\midrule
Monomial coefficients &
  Plain integers: \cd{24u\^{}2} &
  Rational notation: \cd{(24//1)*(u\^{}2)} \\[6pt]
Evaluation table &
  One column per axis (decimal only) &
  Two columns per axis: exact fraction + decimal \\[6pt]
$u$ values &
  Decimals: \cd{0.25}, \cd{0.50} \ldots &
  Exact fractions: \cd{1//4}, \cd{1//2} \ldots \\[6pt]
Verification block &
  Two lines: expression + decimal &
  Three lines: expression, exact fraction, then decimal \\
\bottomrule
\end{tabular}

\medskip
Steps~1, 2, and~3 look almost identical between the two scripts ---
those steps are symbolic descriptions and the expansion results
are the same polynomials.
The real difference is visible in Steps~4 and~5.
\end{diffbox}


% =============================================================
\needspace{8\baselineskip}
\section{Header and Control Points}
% =============================================================

\begin{minted}{text}
==================================================================
  Bezier Curve - Algebraic Expansion  (via Symbolics.jl)
  Case 2  |  Quartic  |  3D  |  5 Control Points
==================================================================

  Control Points:
    P0 :  x = 0  y = 0  z = 1
    P1 :  x = 0  y = 4  z = 1
    P2 :  x = 4  y = 0  z = 1
    P3 :  x = 4  y = 4  z = 1
    P4 :  x = 5  y = 4  z = 1
\end{minted}

The subtitle \textit{via Symbolics.jl} is the only header difference
from Script~01.
The control points are identical.
Internally, however, these coordinates are stored as
\cd{Rational\{Int\}} (exact integer fractions) rather than
\cd{Float64} --- so \cd{4} here means exactly $\frac{4}{1}$,
not the floating-point approximation $4.0$.

This distinction matters later: every coefficient in every
polynomial will be an exact fraction, produced entirely by
integer arithmetic.


% =============================================================
\needspace{8\baselineskip}
\section{Steps 1, 2, and 3 --- Identical in Purpose}
% =============================================================

\begin{minted}{text}
  STEP 1 - Bernstein Form
  ----------------------------------------------------------------
  P(u)  =  [x(u)  y(u)  z(u)]

  P(u)  =  B(4,0)*P0 + B(4,1)*P1 + ... + B(4,4)*P4

  where:
    B(4,0,u)  =  (1-u)^4
    B(4,1,u)  =  4u(1-u)^3
    ...
\end{minted}

\begin{minted}{text}
  STEP 2 - Substitution of Control Point Coordinates
  ----------------------------------------------------------------
  x(u)  =  (1-u)^4*0  +  4u(1-u)^3*0  +  6u^2(1-u)^2*4  + ...
  ...
  (Each basis term is now expanded by Symbolics.jl)
\end{minted}

\begin{minted}{text}
  STEP 3 - Expand Each Bernstein Basis Term
  ----------------------------------------------------------------
  B(4,0,u)  =  (1-u)^4
           =  1 - 4u + 6(u^2) - 4(u^3) + u^4

  B(4,1,u)  =  4u(1-u)^3
           =  4u - 12(u^2) + 12(u^3) - 4(u^4)
  ...
\end{minted}

These three steps are structurally identical to Script~01.
The Bernstein formula, the substitution, and the expansion results
are all the same polynomials.

One small visual difference in Step~3: Symbolics.jl wraps
higher powers in parentheses --- \cd{6(u\^{}2)} rather than \cd{6u\^{}2}.
This is just a display choice by the CAS; mathematically the
expressions are identical.

\begin{notebox}
\textbf{Step 2 closing line:}
\cd{(Each basis term is now expanded by Symbolics.jl)}

In Script~01 the equivalent line said
\cd{(Each basis term will now be expanded individually)},
referring to the hand-written \cd{expand\_bernstein} function.
Here it says \cd{expanded by Symbolics.jl} ---
the CAS calls its own \cd{expand()} function internally,
which applies the binomial theorem and collects like terms
in a single operation.
\end{notebox}


% =============================================================
\needspace{8\baselineskip}
\section{Step 4 --- Collecting Terms}
% =============================================================

\begin{minted}{text}
  STEP 4 - Collect Terms - Final Monomial Polynomial
  ----------------------------------------------------------------
  x(u) - weighted contributions:
    B(4,0)*   0  =  0
    B(4,1)*   0  =  0
    B(4,2)*   4  =  (24//1)*(u^2) - (48//1)*(u^3) + (24//1)*(u^4)
    B(4,3)*   4  =  (16//1)*(u^3) - (16//1)*(u^4)
    B(4,4)*   5  =  (5//1)*(u^4)
\end{minted}

\needspace{5\baselineskip}
\subsection{Understanding \texttt{(24//1)*(u\^{}2)}}

This notation is the most unfamiliar part of the output
for anyone coming from Script~01.
It needs to be read in two pieces.

\textbf{The \cd{//} operator --- Julia's rational number literal.}
In Julia, \cd{a//b} creates an \emph{exact rational number}
$\frac{a}{b}$ stored as a pair of integers.
It never loses precision.

\begin{examplebox}
\cd{24//1} is the rational number $\frac{24}{1} = 24$, exactly.

\cd{1//4} is the rational number $\frac{1}{4} = 0.25$, exactly.
Unlike \cd{Float64(0.25)}, this is represented without any binary
rounding --- it is a genuine fraction in memory.

\cd{1//3} is $\frac{1}{3}$, exactly.
\cd{Float64(1/3)} would be \cd{0.3333333333333333...} with rounding.
\cd{1//3} has no rounding at all.
\end{examplebox}

\textbf{Why does \cd{24} become \cd{24//1}?}
When Symbolics.jl expands an expression whose coefficients
come from a \cd{Rational\{Int\}} matrix, it preserves the
rational type throughout.
The coordinate \cd{4} was stored as \cd{4//1};
multiplying it by the integer coefficient \cd{6}
(from $\binom{4}{2} = 6$) gives \cd{6 * (4//1) = 24//1}.
Symbolics.jl then prints this rational coefficient
as \cd{(24//1)} to make the type explicit.

In short: \cd{(24//1)*(u\^{}2)} means exactly $24u^2$.
The fraction notation is a transparency feature,
not a complication.

\begin{examplebox}
\textbf{Reading Step~4 contributions for $x(u)$:}

\medskip
\begin{tabular}{lll}
\toprule
Symbolic output & Meaning & Script 01 equivalent \\
\midrule
\cd{(24//1)*(u\^{}2) - (48//1)*(u\^{}3) + (24//1)*(u\^{}4)} &
  $24u^2 - 48u^3 + 24u^4$ &
  \cd{24u\^{}2 - 48u\^{}3 + 24u\^{}4} \\[4pt]
\cd{(16//1)*(u\^{}3) - (16//1)*(u\^{}4)} &
  $16u^3 - 16u^4$ &
  \cd{16u\^{}3 - 16u\^{}4} \\[4pt]
\cd{(5//1)*(u\^{}4)} &
  $5u^4$ &
  \cd{5u\^{}4} \\
\bottomrule
\end{tabular}

\medskip
The numbers are identical to Script~01.
Only the display format changed because the type changed.
\end{examplebox}

\needspace{5\baselineskip}
\subsection{Collected Polynomials}

\begin{minted}{text}
  Collected:
    x(u)  =  (24//1)*(u^2) - (32//1)*(u^3) + (13//1)*(u^4)
    y(u)  =  (16//1)*u - (48//1)*(u^2) + (64//1)*(u^3) - (28//1)*(u^4)
    z(u)  =  1
\end{minted}

Same final polynomials as Script~01:
\[
  x(u) = 24u^2 - 32u^3 + 13u^4
\]
\[
  y(u) = 16u - 48u^2 + 64u^3 - 28u^4
\]
\[
  z(u) = 1
\]

The $z(u) = 1$ result is particularly clean ---
Symbolics.jl simply outputs \cd{1} with no rational wrapper,
because the sum of all Bernstein polynomials is identically 1
and the CAS recognises this and simplifies it completely.


% =============================================================
\needspace{8\baselineskip}
\section{Step 5 --- Evaluation Table with Exact Fractions}
% =============================================================

This is the step that most visibly sets Script~02 apart.

\begin{minted}{text}
       u  x(u) exact       decimal  y(u) exact       decimal  z(u) exact   decimal
  ------------------------------------------------------------------------------------
       0           0      0.000000           0      0.000000           1  1.000000
    1//4    269//256      1.050781     121//64      1.890625           1  1.000000
    1//2      45//16      2.812500        9//4      2.250000           1  1.000000
    3//4   1053//256      4.113281     201//64      3.140625           1  1.000000
       1           5      5.000000           4      4.000000           1  1.000000
\end{minted}

\needspace{5\baselineskip}
\subsection{Two columns per axis}

Each axis now has two columns: \textbf{exact} and \textbf{decimal}.

The \textbf{exact} column is the result of substituting the
exact rational $u$ value (e.g.\ \cd{1//4}) into the symbolic
polynomial and simplifying --- the CAS returns an exact fraction.

The \textbf{decimal} column converts that fraction to \cd{Float64}
only at the very last moment, purely for human readability.

\begin{diffbox}
\textbf{This is the core advantage of symbolic computation.}

In Script~01, $x(0.25)$ was computed as:
\[
  24 \times 0.0625 - 32 \times 0.015625 + 13 \times 0.00390625
  = 1.05078125
\]
using floating-point arithmetic at every multiplication.

In Script~02, $x\!\left(\tfrac{1}{4}\right)$ is computed as:
\[
  24 \times \frac{1}{16} - 32 \times \frac{1}{64}
  + 13 \times \frac{1}{256}
  = \frac{384}{256} - \frac{128}{256} + \frac{13}{256}
  = \frac{269}{256}
\]
using exact integer fraction arithmetic.
The decimal $1.050781$ is then read off from $\frac{269}{256}$.
No rounding occurred anywhere in the computation.
\end{diffbox}

\needspace{5\baselineskip}
\subsection{Verifying the exact fractions by hand}

\begin{examplebox}
\textbf{Verify $x\!\left(\tfrac{1}{4}\right) = \dfrac{269}{256}$:}

\medskip
Substitute $u = \tfrac{1}{4}$ into $x(u) = 24u^2 - 32u^3 + 13u^4$:

\[
  24 \cdot \left(\tfrac{1}{4}\right)^2
  - 32 \cdot \left(\tfrac{1}{4}\right)^3
  + 13 \cdot \left(\tfrac{1}{4}\right)^4
  = \frac{24}{16} - \frac{32}{64} + \frac{13}{256}
\]

Converting to a common denominator of 256:
\[
  = \frac{384}{256} - \frac{128}{256} + \frac{13}{256}
  = \frac{384 - 128 + 13}{256}
  = \frac{269}{256}
\]

As decimal: $269 \div 256 = 1.05078125 \approx 1.050781$ \checkmark
\end{examplebox}

\begin{examplebox}
\textbf{Verify $y\!\left(\tfrac{1}{2}\right) = \dfrac{9}{4}$:}

\medskip
Substitute $u = \tfrac{1}{2}$ into $y(u) = 16u - 48u^2 + 64u^3 - 28u^4$:

\[
  16 \cdot \tfrac{1}{2}
  - 48 \cdot \tfrac{1}{4}
  + 64 \cdot \tfrac{1}{8}
  - 28 \cdot \tfrac{1}{16}
  = 8 - 12 + 8 - \tfrac{7}{4}
  = 4 - \tfrac{7}{4}
  = \tfrac{9}{4}
\]

As decimal: $9 \div 4 = 2.25$ \checkmark
\end{examplebox}

\begin{examplebox}
\textbf{Verify $x(1) = 5$ and $y(1) = 4$:}

\[
  x(1) = 24 - 32 + 13 = 5 \checkmark
\]
\[
  y(1) = 16 - 48 + 64 - 28 = 4 \checkmark
\]

At the endpoints the exact values are whole integers ---
no fractions needed. This confirms the curve starts at $P_0 = (0,0,1)$
and ends at $P_4 = (5,4,1)$.
\end{examplebox}

\needspace{5\baselineskip}
\subsection{The $z$ column}

Every entry in the $z$ column is exactly \cd{1} with decimal
\cd{1.000000}.
This is the partition-of-unity property in action:
all five $z$-coordinates are 1, so $z(u) = 1$ for all $u$,
and substituting any rational $u$ into the constant polynomial
$z(u) = 1$ gives exactly \cd{1} --- no fraction, no decimal needed.


% =============================================================
\needspace{8\baselineskip}
\section{Verification Block}
% =============================================================

\begin{minted}{text}
  Verification at u = 1//2:
    x(1//2)  =  (24//1)*(u^2) - (32//1)*(u^3) + (13//1)*(u^4)
             =  45//16
             =  45//16  ~  2.812500

    y(1//2)  =  (16//1)*u - (48//1)*(u^2) + (64//1)*(u^3) - (28//1)*(u^4)
             =  9//4
             =  9//4  ~  2.250000

    z(1//2)  =  1
             =  1
             =  1  ~  1.000000
\end{minted}

The verification block now has \textbf{three lines} per axis,
compared to Script~01's two.

\begin{center}
\begin{tabular}{lp{4cm}p{5cm}}
\toprule
Line & Content & Meaning \\
\midrule
1 & Symbolic expression &
  The full polynomial as Symbolics.jl stored it \\[4pt]
2 & Exact rational result &
  The CAS substituted $u = \tfrac{1}{2}$ and simplified to a fraction \\[4pt]
3 & Fraction $\approx$ decimal &
  The same fraction converted to \cd{Float64} for readability \\
\bottomrule
\end{tabular}
\end{center}

\begin{examplebox}
\textbf{Tracing $x\!\left(\tfrac{1}{2}\right)$ through the three lines:}

\medskip
\textbf{Line 1} --- the stored symbolic expression:
\[
  (24/\!/1)\cdot u^2 - (32/\!/1)\cdot u^3 + (13/\!/1)\cdot u^4
  \quad\equiv\quad 24u^2 - 32u^3 + 13u^4
\]

\textbf{Line 2} --- substitute $u = \tfrac{1}{2}$ and simplify:
\[
  24 \cdot \tfrac{1}{4} - 32 \cdot \tfrac{1}{8} + 13 \cdot \tfrac{1}{16}
  = 6 - 4 + \tfrac{13}{16}
  = 2 + \tfrac{13}{16}
  = \tfrac{45}{16}
\]
So \cd{45//16} is correct. \checkmark

\textbf{Line 3} --- convert to decimal:
$45 \div 16 = 2.8125$, displayed as \cd{2.812500}. \checkmark
\end{examplebox}

\begin{notebox}
\textbf{Why show the fraction twice (lines 2 and 3)?}

Line~2 is the CAS result as returned by \cd{substitute()} ---
it is still a Symbolics expression, printed as \cd{45//16}.
Line~3 is after calling \cd{to\_exact\_rational()} to extract
the plain Julia \cd{Rational\{Int\}} value, and then
\cd{Float64()} for the decimal.
The two-step display confirms that the CAS result and the
extracted rational agree before the decimal conversion.
\end{notebox}


% =============================================================
\needspace{8\baselineskip}
\section{Why Exact Arithmetic Matters}
% =============================================================

\textit{You might wonder: both scripts give the same curve.
So why bother with the extra complexity of symbolic computation?}

The answer is trust and precision.

With Script~01 (floating-point), the decimal \cd{1.050781} could
in principle be a rounded approximation of some irrational number
--- you cannot tell just from the decimal alone.

With Script~02 (exact), the result \cd{269//256} is a \emph{proof}:
the exact value of $x\!\left(\tfrac{1}{4}\right)$ on this curve
is the fraction $\frac{269}{256}$.
No rounding ever occurred.
You can verify it with pencil and paper.

For a quartic curve with integer control points and rational
$u$ values, the output is always a rational number.
Script~02 \emph{finds and displays} that rational number exactly.
Script~01 \emph{approximates} it with a float.

Both are correct for practical use.
But Script~02 gives you something Script~01 cannot:
a result you can algebraically verify without worrying
about floating-point error.

\begin{notebox}
\textbf{A note on the \cd{(24//1)} display.}

You might wish Symbolics.jl printed \cd{24} instead of \cd{(24//1)}.
That is a reasonable aesthetic preference.
The fraction notation is the CAS being explicit about types ---
it is telling you these are rational coefficients,
not floating-point approximations.
Once you have read \cd{24//1} as ``exactly 24'' a few times,
it stops feeling strange.

\textit{Like meeting someone with an unusual name.
After the third time you just call them by it without thinking.}
\end{notebox}

\vspace{2em}\hrule\vspace{0.5em}
{\small\textcolor{gray}{%
Output explanation for \texttt{02-bezier-sympy.jl}, Case~2
\textbar{} Code by E.R.\ Nurwijayadi
\textbar{} Explanations prompt-directed by E.R.\ Nurwijayadi, AI-composed}}

\end{document}
